\chapter{Theoretical Background}\label{ch:theoretical-background}

\section{State of the Art}
% todo add:
% History of prediction (competitions, etc.)
% history of forecasting
% bring together to introduce Müller & Yangs paper
% description of EEG (also subscalp, scalp, etc.)


\subsection{Evolution of Seizure Prediction and Algorithmic Development}

The pursuit of extracting seizure precursors from the electroencephalogram began in the 1970s with \textit{linear approaches}, such as autoregressive modeling and power-spectrum analysis of spike-wave bursts \parencite{mormann_seizure_2007}.
Early research also focused on the predictive value of \textit{interictal spike rates}, though results were often contradictory; while some early studies suggested focal spiking decreased before seizures, subsequent investigations on more extended databases showed no systematic or reproducible changes \parencite{mormann_seizure_2007}.

In the 1990s, the field shifted toward \textit{nonlinear dynamics}, treating the brain as a chaotic system \parencite{mormann_seizure_2007, kuhlmann_epilepsyecosystemorg_2018}.
Researchers utilized measures like the largest Lyapunov exponent to identify a decrease in chaoticity in the minutes preceding a seizure, alongside the dynamical similarity index and correlation density to characterize a transitional preictal phase \parencite{mormann_seizure_2007}.
This era was marked by significant optimism, with many studies reporting high sensitivity on small, highly selected datasets \parencite{mormann_seizure_2007, muller_coherent_2022}.
However, a period of \textit{skepticism} followed in the early 2000s as more rigorous evaluations on unselected, continuous multi-day recordings failed to reproduce these optimistic findings, highlighting issues with overfitting, lack of specificity evaluation, and the absence of out-of-sample statistical validation \parencite{mormann_seizure_2007, freestone_forward-looking_2017}.

To overcome these reproducibility challenges and address the \enquote{ceiling effect} in predictive performance, the field turned to \textit{crowdsourcing and open competitions} using standardized, high-quality datasets. Platforms like \textit{Kaggle} hosted the American Epilepsy Society Seizure Prediction Challenge, providing \gls{ieeg} recordings from both humans and canines \parencite{brinkmann_crowdsourcing_2016}.
Subsequent efforts by \textit{epilepsyecosystem.org} continued this trend, enabling researchers to compare algorithms on ultra-long-term human recordings \parencite{kuhlmann_epilepsyecosystemorg_2018}.
These competitions demonstrated that \textit{machine learning} could achieve better-than-chance performance \parencite{brinkmann_crowdsourcing_2016}.
Winning solutions typically utilized a variety of \enquote{hand-crafted} features---including band power spectra, statistical moments (skewness, kurtosis), and cross-channel linear correlations in both time and frequency domains---classified by \textit{Support Vector Machines (SVM)}, \textit{Random Forests}, or \textit{k-nearest neighbors} \parencite{kuhlmann_epilepsyecosystemorg_2018, muller_coherent_2022}.

More recently, \textit{deep learning} has emerged as a dominant paradigm, with \textit{\gls{cnn}} and \textit{\gls{lstm}} networks used to extract features directly from raw data, bypassing the need for manual feature engineering \parencite{kuhlmann_epilepsyecosystemorg_2018, muller_coherent_2022}.
Looking forward, \textcite{freestone_forward-looking_2017} suggest the field is shifting from binary classification to a \textit{probabilistic framework}, casting prediction as a regression problem that expresses a continuous seizure likelihood.
This future vision relies on power-efficient, \textit{neuromorphic hardware}---such as IBM's TrueNorth chip---and less-invasive \textit{subscalp monitors} to support sophisticated real-time monitoring in ambulatory settings \parencite{freestone_forward-looking_2017, yang_seizure_2024}.


\subsection{Precursors and Mechanisms of Ictogenesis}
% todo keep going through section and fix errors
Understanding the mechanisms by which a \enquote{normal} brain transitions into a seizure---a process known as \textit{ictogenesis}---is fundamental to predictive modeling \parencite{mormann_seizure_2007}.
While some patients report clinical prodromi such as behavioral or mood changes, the field operates on the assumption that these transitions are best tracked by \textit{\gls{eeg}}, which can capture measurable shifts in brain dynamics hours or even days before a clinical event in many patients \parencite{freestone_forward-looking_2017, cook_prediction_2013}.

On a network level, focal seizures are hypothesized to arise from \textit{\enquote{burster} neurons} that abnormally recruit and entrain neighbors into a \enquote{critical mass}.
This build-up is viewed not merely as an increasing count of neurons but as an unfolding process of increasing \textit{critical interactions} within a pathological functional network \parencite{mormann_seizure_2007}.

A key indicator of this transition is \enquote{\textit{critical slowing}}, a universal signature of a system losing its resilience as it approaches a tipping point or state transition.
In the context of epilepsy, this is often characterized by an increase in the autocorrelation of the \gls{eeg} or a lengthening of electrical evoked potentials, indicating a system that takes longer to recover from perturbations  \parencite{freestone_forward-looking_2017, maturana_critical_2020}.

This transition is frequently framed through the lens of \textit{bi/multistability}, where a stable healthy attractor (e.g., wakefulness or sleep) coexists with a pathological seizure state.
Seizures may occur due to deterministic variations in underlying parameters or random \enquote{noisy} fluctuations in background activity that drive the brain's state across the energy barrier separating these attractors.
Because the brain possesses internal \textit{homeostatic regulatory mechanisms}, it is possible for the brain to enter a state of \textit{high susceptibility} without a clinical event actually occurring, as these mechanisms may correct the abnormal activity \parencite{freestone_forward-looking_2017}.

This conceptualization leads to the idea of a \textit{proictal state}---a period of heightened seizure likelihood that does not always terminate in an ictal event \parencite{muller_coherent_2022}.
Furthermore, research into biomarkers like \textit{interictal spikes} and \textit{\glspl{hfo}} has revealed that while these indicators have clear implications for understanding ictogenesis, their dynamics are highly patient-specific \parencite{freestone_forward-looking_2017}.

Consequently, there is an increasing recognition that \enquote{one-size-fits-all} algorithms are insufficient, necessitating long-term monitoring to account for individual variability in seizure-generating mechanisms and the diverse populations of seizures that can coexist within a single patient \parencite{freestone_forward-looking_2017, yang_seizure_2024}.



\section{Electroencephalography (EEG)}\label{sec:eeg-background}

\textbf{\Gls{eeg}} is an electro-physiological technique used to record the electrical activity arising from brains. 
It primarily measures the \textbf{summation of excitatory and inhibitory postsynaptic potentials} from relatively large groups of neurons---especially cortical pyramidal neurons---firing synchronously. 
Because these neurons are oriented perpendicularly to the brain's surface, their combined activity creates a dipole that can be detected by recording electrodes. 
\gls{eeg} is characterized by its \textbf{exquisite temporal sensitivity}, providing a near real-time display of ongoing cerebral activity \parencite{st_louis_electroencephalography_2016}.

\subsection{Clinical Utility and Applications}
The primary utility of \gls{eeg} lies in the evaluation of \textbf{dynamic cerebral functioning}. 
It is widely used for several clinical indications \parencite{st_louis_electroencephalography_2016}:
\begin{itemize}
    \item \textbf{Epilepsy:} \gls{eeg} is the gold standard for diagnosing and confirming seizures and epilepsy. 
    It allows clinicians to identify characteristic ictal (during-seizure) alterations and \glspl{ied}, such as spikes and sharp waves, which help classify seizure types and epilepsy syndromes.
    \item \textbf{General monitoring:} It is used to monitor the \textit{depth of anesthesia} during surgery and to detect sudden changes in neural functioning that may indicate \textit{ischemia or infarction}.
    \item \textbf{Research and specialized diagnosis:} \gls{eeg} waveforms can be averaged to study \textit{\glspl{ep}} and \textit{\glspl{erp}}, which represent neural activity related to specific stimuli, assisting in the analysis of cognitive and sensory functioning.
\end{itemize}

\subsection{Frequency Bands}
\gls{eeg} activity is conventionally categorized into frequency\footnote{The exact frequencies boundaries of each band vary in the literature and are based on \textcite{mormann_seizure_2007} here.} bandwidths, or bands, which vary according to the patient's state (e.g., wakefulness, sleep) and developmental stage \parencite{st_louis_electroencephalography_2016,, mormann_seizure_2007, nayak_normal_2026}:

{
\sisetup{range-phrase = {--}, range-units = single}
\begin{itemize}
    \item \textbf{Delta $\delta$} (\SIrange{0.5}{4}{\hertz}): Predominant during deep sleep and considered pathological if present in awake adults.
    \item \textbf{Theta $\theta$} (\SIrange{4}{8}{\hertz}): Commonly seen during drowsiness and in early sleep stages.
    \item \textbf{Alpha $\alpha$} (\SIrange{8}{12}{\hertz}): Present during relaxed, awake states, especially in the occipital head region. Strengthens when eyes are closed and attenuates during intense activity or concentration.
    \item \textbf{Beta $\beta$} (\SIrange{12}{35}{\hertz}): Associated with alert wakefulness and often enhanced by most sedative medications.
    \item \textbf{Gamma $\gamma$} (\SIrange{35}{100}{\hertz}): Associated with active cognition. Is not commonly seen in scalp \gls{eeg} but is readily detected in intracranial recordings.
\end{itemize}
}

\subsection{EEG Recording Modalities}
The choice of electrode placement significantly impacts the recorded signal quality and spatial resolution \parencite{st_louis_electroencephalography_2016}:
\begin{itemize}
    \item \textbf{Surface (scalp) \gls{eeg}:} Non-invasive electrodes are placed on the skin using systems like the \textit{International 10--20 System}. 
    However, the signal must traverse biological filters (brain, cerebrospinal fluid, meninges, skull, and skin), which reduce its amplitude and blur spatial information. 
    The skull acts as a \textit{spatial averager} and a dominant attenuator of high-frequency signals.
    \item \textbf{\Gls{ecog} / \gls{ieeg}:} Invasive monitoring involves placing electrodes directly on the cortex (subdural grids or strips) or within the brain (depth electrodes). 
    These provide much cleaner fidelity without muscle or movement artifacts, enabling precise localization of epileptic foci.
    \item \textbf{Subcutaneous (subscalp) \gls{eeg}:} This modality involves electrodes implanted under the skin but outside the skull, such as the \textit{UNEEG SubQ system} used in this study.
    Such systems offering a robust and unobtrusive alternative to traditional scalp \gls{eeg} \parencite{duun-henriksen_subdural_2013}. 
\end{itemize}

Research comparing subcutaneous systems to standard surface electrodes shows that they achieve similar performance, as measured by \textit{\gls{psd}} and \textit{coherence analysis}. 
Subcutaneous recordings often demonstrate \textit{higher signal quality} during rest, primarily because the electrodes provide \textit{higher signal amplitudes} and maintain a \textit{rigid connection} to the source, reducing artifacts caused by electrode-skin instability or movements \parencite{duun-henriksen_eeg_2015}. 
While the skull still acts as a spatial averager for these signals, the subcutaneous approach avoids the degradation of the electrodes contact to the of skin over time, common with surface electrodes, as well as being portable and inconspicuous, making it suitable for \enquote{ultra-long-term} monitoring over months.
